% Chunk: Physical pages 465--470, beginning with the continuation of Section 10.3.
% Source: Weinberg QFT Vol. I, physical PDF pages 465--470 (printed pages 442--447).
% Figures: Figure 10.5, with caption transcribed and labeled.
% Uncertainties: none.

\begin{equation}
\mathcal L_1=-(Z_2-1)[\overline\Psi[\sl{\partial}+m]\Psi]
+Z_2\delta m\overline\Psi\Psi
-V_{\mathrm B}(Z_2\overline\Psi\Psi).
\tag{10.3.27}\label{eq:10.3.27}
\end{equation}
Let $i(2\pi)^4\Sigma^*(\sl{k})$ be the sum of all connected graphs, with one
fermion line coming in with four-momentum $k$ and one going out with the same
four-momentum, that cannot be disconnected by cutting through any single
internal fermion line, and with external line propagator factors
$-i(2\pi)^{-4}$ and $[i\sl{k}+m-i\epsilon]^{-1}$ omitted. (Lorentz invariance
is being used to justify writing $\Sigma^*$ as an ordinary function of the
Lorentz scalar matrix $\sl{k}=k_\mu\gamma^\mu$.) Then the complete fermion
propagator is
\begin{equation}
\begin{split}
S'(k)={}&[i\sl{k}+m-i\epsilon]^{-1}
+[i\sl{k}+m-i\epsilon]^{-1}\Sigma^*(\sl{k})
[i\sl{k}+m-i\epsilon]^{-1}
\\
&+[i\sl{k}+m-i\epsilon]^{-1}\Sigma^*(\sl{k})
[i\sl{k}+m-i\epsilon]^{-1}\Sigma^*(\sl{k})
[i\sl{k}+m-i\epsilon]^{-1}+\cdots
\\
={}&[i\sl{k}+m-\Sigma^*(\sl{k})-i\epsilon]^{-1}.
\end{split}
\tag{10.3.28}\label{eq:10.3.28}
\end{equation}
In calculating $\Sigma^*(\sl{k})$ we take into account the tree graphs from
the terms in Eq.~\eqref{eq:10.3.27} proportional to
$\overline\Psi\sl{\partial}\Psi$ and $\overline\Psi\Psi$ as well as loop
contributions:
\begin{equation}
\Sigma^*(\sl{k})=-(Z_2-1)[i\sl{k}+m]+Z_2\delta m
+\Sigma^*_{\mathrm{LOOP}}(\sl{k}).
\tag{10.3.29}\label{eq:10.3.29}
\end{equation}
The condition that the complete propagator has a pole at $k^2=-m^2$ with the
same residue as the uncorrected propagator is then that
\begin{equation}
\Sigma^*(im)=0,
\tag{10.3.30}\label{eq:10.3.30}
\end{equation}
\begin{equation}
\left.\frac{\partial\Sigma^*(\sl{k})}{\partial\sl{k}}
\right|_{\sl{k}=im}=0,
\tag{10.3.31}\label{eq:10.3.31}
\end{equation}
and hence
\begin{equation}
Z_2\delta m=-\Sigma^*_{\mathrm{LOOP}}(im),
\tag{10.3.32}\label{eq:10.3.32}
\end{equation}
\begin{equation}
Z_2=1-i\left.\frac{\partial\Sigma^*_{\mathrm{LOOP}}(\sl{k})}
{\partial\sl{k}}\right|_{\sl{k}=im}.
\tag{10.3.33}\label{eq:10.3.33}
\end{equation}
Just as for scalars, the vanishing of
$[i\sl{k}+m]^{-1}\Sigma^*(\sl{k})$ in the limit $\sl{k}\to im$ tells us that
radiative corrections may be ignored in external fermion lines. Corresponding
results for the photon propagator will be derived in Section 10.5.

\subsection{Renormalized Charge and Ward Identities}

The use of the commutation and conservation relations of Heisenberg-picture
operators allows us to make a connection between the charges (or other similar
quantities) in the Lagrangian density and the properties of physical states.
Recall that the invariance of the Lagrangian density with respect to global
gauge transformations $\Psi_\ell\to\exp(iq_\ell\alpha)\Psi_\ell$ (with
$\alpha$ an arbitrary constant phase) implies the existence of a current
\begin{equation}
J^\mu=-i\sum_\ell
\frac{\partial\mathcal L}{\partial(\partial_\mu\Psi_\ell)}
q_\ell\Psi_\ell,
\tag{10.4.1}\label{eq:10.4.1}
\end{equation}
satisfying the conservation condition
\begin{equation}
\partial_\mu J^\mu=0.
\tag{10.4.2}\label{eq:10.4.2}
\end{equation}
This implies that the space-integral of the time component of $J^\mu$ is
time-independent:
\begin{equation}
i\frac{d}{dt}Q=[Q,H]=0,
\tag{10.4.3}\label{eq:10.4.3}
\end{equation}
where
\begin{equation}
Q\equiv\int d^3x\,J^0.
\tag{10.4.4}\label{eq:10.4.4}
\end{equation}
(There is a very important possible exception here, that the integral
\eqref{eq:10.4.4} may not exist if there are long-range forces due to massless
scalars in the system. We will return to this point when we consider broken
symmetries in Volume II.) Also, since it is a space-integral, $Q$ is manifestly
translation-invariant
\begin{equation}
[\mathbf P,Q]=0
\tag{10.4.5}\label{eq:10.4.5}
\end{equation}
and since $J^\mu$ is a four-vector, $Q$ is invariant with respect to
homogeneous Lorentz transformations
\begin{equation}
[J^{\mu\nu},Q]=0.
\tag{10.4.6}\label{eq:10.4.6}
\end{equation}
It follows that $Q$ acting on the true vacuum $\ket{\Omega}$ must be another
Lorentz-invariant state of zero energy and momentum, and hence (assuming no
vacuum degeneracy) must be proportional to $\ket{\Omega}$ itself. But the
proportionality constant must vanish, because Lorentz invariance requires that
$\bra{\Omega}J_\mu\ket{\Omega}$ vanish. Hence
\begin{equation}
Q\ket{\Omega}=0.
\tag{10.4.7}\label{eq:10.4.7}
\end{equation}
Also, $Q$ acting on any one-particle state $\ket{\mathbf p,\sigma,n}$ must be
another state with the same energy, momentum, and Lorentz transformation
properties, and thus (assuming no degeneracy of one-particle states) must be
proportional to the same one-particle state
\begin{equation}
Q\ket{\mathbf p,\sigma,n}=q_{(n)}\ket{\mathbf p,\sigma,n}.
\tag{10.4.8}\label{eq:10.4.8}
\end{equation}
The Lorentz invariance of $Q$ ensures that the eigenvalue $q_{(n)}$ is
independent of $\mathbf p$ and $\sigma$, depending only on the species of the
particle. This eigenvalue is what is known as the electric charge (or whatever
other quantum number of which $J^\mu$ may be the current) of the one-particle
state. To relate this to the $q_\ell$ parameters in the Lagrangian, we note that
the canonical commutation relations give
\begin{equation}
[J^0(\mathbf x,t),\Psi_\ell(\mathbf y,t)]
=-q_\ell\Psi_\ell(\mathbf y,t)\delta^3(\mathbf x-\mathbf y),
\tag{10.4.9}\label{eq:10.4.9}
\end{equation}
or integrating over $\mathbf x$:
\begin{equation}
[Q,\Psi_\ell(y)]=-q_\ell\Psi_\ell(y).
\tag{10.4.10}\label{eq:10.4.10}
\end{equation}
The same is true of any local function $F(y)$ of the fields and field
derivatives and their adjoints, containing definite numbers of each:
\begin{equation}
[Q,F(y)]=-q_FF(y),
\tag{10.4.11}\label{eq:10.4.11}
\end{equation}
where $q_F$ is the sum of the $q_\ell$ for all fields and field derivatives in
$F(y)$, minus the sum of the $q_\ell$ for all field adjoints and their
derivatives. Taking the matrix element of this equation between a one-particle
state and the vacuum, and using Eqs.~\eqref{eq:10.4.7} and
\eqref{eq:10.4.8}, we have
\begin{equation}
\bra{\Omega}F(y)\ket{\mathbf p,\sigma,n}(q_F-q_{(n)})=0.
\tag{10.4.12}\label{eq:10.4.12}
\end{equation}
Hence we must have
\begin{equation}
q_{(n)}=q_F
\tag{10.4.13}\label{eq:10.4.13}
\end{equation}
as long as
\begin{equation}
\bra{\Omega}F(y)\ket{\mathbf p,\sigma,n}\ne0.
\tag{10.4.14}\label{eq:10.4.14}
\end{equation}
As we saw in the previous section, Eq.~\eqref{eq:10.4.14} is the condition
that assures that momentum space Greens functions involving $F$ have poles
corresponding to the one-particle state $\ket{\mathbf p,\sigma,n}$. For a
one-particle state corresponding to one of the fields in the Lagrangian we
could take $F=\Psi_\ell$, in which case $q_F=q_\ell$, but our results here
apply to general one-particle states, whether or not their fields appear in the
Lagrangian.

This almost, but not quite, tells us that despite all the possible high-order
graphs that affect the emission and absorption of photons by charged particles,
the physical electric charge is just equal to a parameter $q_\ell$ appearing in
the Lagrangian (or to a sum of such parameters, like $q_F$.) The qualification
that has to be added here is that the requirement, that the Lagrangian be
invariant under the transformations
$\Psi_\ell\to\exp(iq_\ell\alpha)\Psi_\ell$, does nothing to fix the over-all
scale of the quantities $q_\ell$. The physical electric charges are those that
determine the response of matter fields to a given \emph{renormalized}
electromagnetic field $A^\mu$. That is, the scale of the $q_\ell$ is fixed by
requiring that the renormalized electromagnetic field appears in the matter
Lagrangian $\mathcal L_M$ in the linear combinations
$[\partial_\mu-iq_\ell A_\mu]\Psi_\ell$, so that the current $J^\mu$ is
\begin{equation}
J^\mu=\frac{\delta\mathcal L_M}{\delta A_\mu}.
\tag{10.4.15}\label{eq:10.4.15}
\end{equation}
But $A^\mu$ and $q_\ell$ are not the same as the `bare electromagnetic field'
$A_{\mathrm B}^\mu$ and `bare charges' $q_{\mathrm B\ell}$ that appear in the
Lagrangian when we write it in its simplest form
\begin{equation}
\begin{split}
\mathcal L={}&-\frac14(\partial_\mu A_{{\mathrm B}\nu}
-\partial_\nu A_{{\mathrm B}\mu})
(\partial^\mu A_{\mathrm B}^{\nu}-\partial^\nu A_{\mathrm B}^{\mu})
\\
&+\mathcal L_M\bigl(\Psi_\ell,
[\partial_\mu-iq_{\mathrm B\ell}A_{{\mathrm B}\mu}]\Psi_\ell\bigr).
\end{split}
\tag{10.4.16}\label{eq:10.4.16}
\end{equation}
The renormalized electromagnetic field (defined to have a complete propagator
whose pole at $p^2=0$ has unit residue) is conventionally written in terms of
$A_{\mathrm B}^\mu$ as
\begin{equation}
A^\mu=Z_3^{-1/2}A_{\mathrm B}^\mu,
\tag{10.4.17}\label{eq:10.4.17}
\end{equation}
so in order for the charge $q_\ell$ to characterize the response of the
charged particles to a given renormalized electromagnetic field, we should
define the renormalized charges by
\begin{equation}
q_\ell=\sqrt{Z_3}\,q_{\mathrm B\ell}.
\tag{10.4.18}\label{eq:10.4.18}
\end{equation}

We see that the physical electric charge $q$ of any particle is just
proportional to a parameter $q_{\mathrm B}$ related to those appearing in the
Lagrangian, with a proportionality constant $Z_3^{-1/2}$ that is the same for
all particles. This helps us to understand how a particle like the proton,
that is surrounded by a cloud of virtual mesons and other strongly interacting
particles, can have the same charge as the positron, whose interactions are all
much weaker. It is only necessary to assume that for some reason the charges
$q_{\mathrm B\ell}$ in the Lagrangian are equal and opposite for the electron
and for those particles (two $u$ quarks and one $d$ quark) that make up the
proton; the effect of higher-order corrections then appears solely in the
common factor $Z_3^{-1/2}$.

In order for charge renormalization to arise only from radiative corrections
to the photon propagator, there must be cancellations among the great variety
of other radiative corrections to the propagators and electromagnetic vertices
of the charged particles. We can see a little more deeply into the nature of
these cancellations by making use of the celebrated relations between these
charged particle propagators and vertices known as the \emph{Ward identities}.

For instance, consider the Greens function for an electric current $J^\mu(x)$
together with a Heisenberg-picture Dirac field $\Psi_n(y)$ of charge $q$ and
its covariant adjoint $\overline\Psi_m(z)$. We define the electromagnetic
vertex function $\Gamma^\mu$

% Figure 10.5. Source: physical PDF p. 469 (printed p. 446).
\begingroup
\renewcommand{\thefigure}{10.\arabic{figure}}
\setcounter{figure}{4}
\begin{figure}[H]
\centering
\begin{tikzpicture}[
  x=1cm,
  y=1cm,
  line cap=round,
  line join=round,
  electron/.style={draw=black, line width=0.8pt},
  photon/.style={draw=black, line width=0.75pt,
    decorate, decoration={snake, amplitude=1.15mm, segment length=4.4mm}},
  fermion arrow/.style={
    postaction={decorate},
    decoration={markings,
      mark=at position 0.52 with
        {\arrow{Stealth[length=2.4mm,width=1.7mm]}}}
  }
]
  % First correction to the electron propagator: a virtual photon joins
  % two points on the electron line.  Fermion flow is from left to right.
  \coordinate (aL) at (-1.55,1.65);
  \coordinate (aR) at ( 1.55,1.65);
  \draw[electron] (-3.25,1.65) -- (3.25,1.65);
  \draw[electron,fermion arrow] (aL) -- (aR);
  \draw[photon] (aL)
    .. controls (-1.32,2.88) and (1.32,2.88) .. (aR);

  % First correction to the electromagnetic vertex: the external photon
  % enters at the central vertex, between two rightward electron segments.
  \coordinate (bL) at (-1.55,-0.95);
  \coordinate (bC) at ( 0.00,-0.95);
  \coordinate (bR) at ( 1.55,-0.95);
  \draw[electron] (-3.25,-0.95) -- (3.25,-0.95);
  \draw[electron,fermion arrow] (bL) -- (bC);
  \draw[electron,fermion arrow] (bC) -- (bR);
  \draw[photon] (bL)
    .. controls (-1.32,0.28) and (1.32,0.28) .. (bR);
  \draw[photon] (bC) -- (0,-3.10);
\end{tikzpicture}
\caption{Diagrams for the first corrections to the electron propagator and vertex function in quantum electrodynamics. Here straight lines are electrons; wavy lines are photons.}
\label{fig:10.5}
\end{figure}
\endgroup


of the charged particle by

\begin{equation}
\begin{split}
&\int d^4x\,d^4y\,d^4z\,e^{-ip\cdot x}e^{-ik\cdot y}e^{+i\ell\cdot z}
\bra{\Omega}T\{J^\mu(x)\Psi_n(y)\overline\Psi_m(z)\}\ket{\Omega}
\\
&\qquad\equiv-iqS'_{nn'}(k)\Gamma^\mu_{n'm'}(k,\ell)
S'_{m'm}(\ell)\delta^4(p+k-\ell),
\end{split}
\tag{10.4.19}\label{eq:10.4.19}
\end{equation}
where
\begin{equation}
-iS'_{nm}(k)\delta^4(k-\ell)\equiv
\int d^4y\,d^4z\,
\bra{\Omega}T\{\Psi_n(y)\overline\Psi_m(z)\}\ket{\Omega}
e^{-ik\cdot y}e^{+i\ell\cdot z}.
\tag{10.4.20}\label{eq:10.4.20}
\end{equation}
According to the theorem of Section 6.4, Eq.~\eqref{eq:10.4.20} gives the sum
of all Feynman graphs with one incoming and one outgoing fermion line, i.e.,
the complete Dirac propagator. Also, Eq.~\eqref{eq:10.4.19} gives the sum of
all such graphs with an extra photon line attached, so $\Gamma^\mu$ is the sum
of ``vertex'' graphs with one incoming Dirac line, one outgoing Dirac line,
and one photon line, but with the complete Dirac external line propagators and
the bare photon external line propagator stripped away. To make the
normalization of $S'$ and $\Gamma^\mu$ perfectly clear, we mention that in the
limit of no interactions, these functions take the values
\[
S'(k)\to[i\gamma_\nu k^\nu+m-i\epsilon]^{-1},
\qquad
\Gamma^\mu(k,\ell)\to\gamma^\mu.
\]
The one-loop diagrams that provide corrections to these limiting values are
shown in Figure~\ref{fig:10.5}.

We can derive a relation between $\Gamma^\mu$ and $S'$ by use of the identity
\begin{equation}
\begin{split}
\frac{\partial}{\partial x^\mu}
T\{J^\mu(x)\Psi_n(y)\overline\Psi_m(z)\}={}&
T\{\partial_\mu J^\mu(x)\Psi_n(y)\overline\Psi_m(z)\}
\\
&+\delta(x^0-y^0)T\{[J^0(x),\Psi_n(y)]\overline\Psi_m(z)\}
\\
&+\delta(x^0-z^0)T\{\Psi_n(y)[J^0(x),\overline\Psi_m(z)]\}.
\end{split}
\tag{10.4.21}\label{eq:10.4.21}
\end{equation}
where the delta functions arise from time-derivatives of step functions. The
conservation condition \eqref{eq:10.4.2} tells us that the first term vanishes,
while the second and third terms can be calculated using the commutation
relations \eqref{eq:10.4.9}, which here give
\begin{equation}
[J^0(\mathbf x,t),\Psi_n(\mathbf y,t)]
=-q\Psi_n(\mathbf y,t)\delta^3(\mathbf x-\mathbf y)
\tag{10.4.22}\label{eq:10.4.22}
\end{equation}
and its adjoint
\begin{equation}
[J^0(\mathbf x,t),\overline\Psi_n(\mathbf y,t)]
=q\overline\Psi_n(\mathbf y,t)\delta^3(\mathbf x-\mathbf y).
\tag{10.4.23}\label{eq:10.4.23}
\end{equation}
Eq.~\eqref{eq:10.4.21} then reads
\begin{equation}
\begin{split}
\frac{\partial}{\partial x^\mu}
T\{J^\mu(x)\Psi_n(y)\overline\Psi_m(z)\}={}&
-q\delta^4(x-y)T\{\Psi_n(y)\overline\Psi_m(z)\}
\\
&+q\delta^4(x-z)T\{\Psi_n(y)\overline\Psi_m(z)\}.
\end{split}
\tag{10.4.24}\label{eq:10.4.24}
\end{equation}
Inserting this in the Fourier transform \eqref{eq:10.4.19} gives
\[
(\ell-k)_\mu S'(k)\Gamma^\mu(k,\ell)S'(\ell)
=iS'(\ell)-iS'(k)
\]
or in other words
\begin{equation}
(\ell-k)_\mu\Gamma^\mu(k,\ell)
=iS'^{-1}(k)-iS'^{-1}(\ell).
\tag{10.4.25}\label{eq:10.4.25}
\end{equation}
This is known as the \emph{generalized Ward identity}, first derived (by these
methods) by Takahashi~\hyperref[ref:10.4]{[4]}. The original Ward identity,
derived earlier by Ward~\hyperref[ref:10.5]{[5]} from a study of perturbation
theory, can be obtained from Eq.~\eqref{eq:10.4.25} by letting $\ell$ approach
$k$. In this limit, Eq.~\eqref{eq:10.4.25} gives
\begin{equation}
\Gamma^\mu(k,k)=-i\frac{\partial}{\partial k_\mu}S'^{-1}(k).
\tag{10.4.26}\label{eq:10.4.26}
\end{equation}
The fermion propagator is related to the self-energy insertion $\Sigma^*(k)$
by Eq.~\eqref{eq:10.3.28}
\[
S'^{-1}(k)=i\sl{k}+m-\Sigma^*(k),
\]
so Eq.~\eqref{eq:10.4.26} may be written
\begin{equation}
\Gamma^\mu(k,k)=\gamma^\mu+i\frac{\partial}{\partial k_\mu}\Sigma^*(k).
\tag{10.4.27}\label{eq:10.4.27}
\end{equation}
